% !TeX root = ../alda_g17.tex

\pagestyle{short}
\tableofcontents
\clearpage
\pagestyle{long}

\chapter{Introducción}

\section{Descripción del problema}
La captura de movimiento markerless (sin marcadores físicos) representa una línea activa de investigación en visión artificial con aplicaciones en robótica, rehabilitación médica, interfaces humano-máquina e interacción natural. \cite{Pauzietal2021}

Los enfoques tradicionales basados en sensores físicos, trajes de captura óptica o cámaras de profundidad resultan costosos e intrusivos. El presente proyecto propone un sistema basado en una cámara convencional que procesa video en tiempo real, detecta el hombro, codo y muñeca derechos, estima la apertura de la mano y transfiere esos valores al control de un UR5 con gripper RG2 dentro de CoppeliaSim. \cite{Vintimilla2026}

La solución combina OpenCV para la adquisición y el procesamiento de imagen,
MediaPipe Pose y Hands para la percepción, y la ZeroMQ Remote API para controlar
el simulador. Un dashboard integra la cámara humana, la vista del robot, las
mediciones y el estado de seguridad. \cite{Bazarevskyetal2020}
\cite{Lugaresietal2019}

\subsection*{Diagrama de Bloques General del Sistema}
La figura \ref{fig:diagrama_bloques} representa el flujo de datos global del proyecto, descompuesto en sus cinco módulos funcionales principales.

\begin{figure}[h!]
    \centering
    \includegraphics[width=1\textwidth]{images/diagrama_bloques.jpeg}
    \caption{Diagrama de bloques del sistema de imitación de movimientos con MediaPipe Pose.}
    \label{fig:diagrama_bloques}
\end{figure}

Cada bloque representa un módulo funcional independiente con entradas y salidas definidas. El video de la cámara entra al sistema por M1, los landmarks son extraídos en M2, los ángulos se calculan en M3, el brazo robótico es simulado en M4 y todo se integra visualmente en M5. \cite{Bazarevskyetal2020} \cite{Bradski2000}
\section{Justificación del problema}
MediaPipe fue presentado originalmente como un framework modular para construir pipelines de percepción reproducibles entre plataformas. Su componente BlazePose amplió esta capacidad al dominio de la estimación de pose corporal en tiempo real, logrando latencias menores a 30 ms en CPU mediante un pipeline de dos etapas: detector de persona (BlazePose Detector) y estimador de pose (BlazePose Estimator) que extrae 33 puntos clave. \cite{Bazarevskyetal2020} \cite{Lugaresietal2019}

Trabajos previos como los de Pauzi et al. (2021) demostraron la efectividad de BlazePose para estimar movimientos corporales con una diferencia de precisión dentro del 10\% respecto a sistemas IMU de referencia. En el dominio de la teleoperación robótica, Altayeb (2023) implementó un sistema de replicación de gestos de mano en un brazo robótico usando MediaPipe Hands, validando la viabilidad de este enfoque para control de actuadores virtuales. \cite{Pauzietal2021} \cite{Altayeb2023}

La justificación económica radica en el bajo costo de implementación: el sistema requiere únicamente una webcam USB estándar ($<$ \$20 USD), sin sensores de profundidad ni hardware especializado. El stack de software es completamente open-source (Python, OpenCV, MediaPipe, NumPy).

\section{Objetivos}

\subsection{Objetivo principal}
Desarrollar un sistema de visión artificial que, a partir de video capturado en
tiempo real por una cámara convencional, detecte los movimientos del brazo
derecho y la apertura de la mano con MediaPipe, y los transfiera al control
sincronizado y supervisado de un UR5 con gripper RG2 dentro de CoppeliaSim.

\subsection{Objetivos específicos}
\begin{enumerate}
    \item Implementar la captura y el preprocesamiento de video en tiempo real
    con OpenCV, normalización a $640\times480$ y conversión BGR$\rightarrow$RGB.
    \item Integrar MediaPipe Pose y Hands para detectar el hombro, codo y muñeca
    derechos, y asociar la mano correspondiente al brazo observado.
    \item Calcular los ángulos de hombro y codo y la apertura normalizada del
    gripper, con validación geométrica y media móvil de cinco muestras.
    \item Controlar el UR5 y el RG2 en CoppeliaSim mediante ZeroMQ, respetando
    límites articulares e incorporando estados, watchdog, colisiones y ESTOP.
    \item Construir una visualización integrada con cámara humana, Vision
    Sensor, ángulos, apertura, FPS, diagnósticos y acciones de seguridad.
    \item Verificar los contratos del software y la integración con el
    simulador, y definir un protocolo para medir FPS, latencia, MAE y robustez
    bajo condiciones variables.
\end{enumerate}

\section{Módulos del proyecto}
El sistema se descompone en cinco módulos funcionales interconectados. La arquitectura modular facilita el desarrollo y prueba independiente de cada componente en el marco SCRUM adoptado. Las tecnologías principales son OpenCV \cite{Bradski2000} para procesamiento de imagen y MediaPipe Pose \cite{Bazarevskyetal2020, 10.1007/978-3-319-10602-1_48} para estimación de pose.

\subsection{Módulo 1 -- Captura y Preprocesamiento de Imagen}

\subsubsection*{Objetivo}
Adquirir el flujo de video en tiempo real desde la cámara y preparar
cada frame para su procesamiento posterior mediante normalización de
resolución, conversión de espacio de color y reducción de ruido.

\paragraph{Entradas y Salidas}

\begin{table}[h!]
\centering
\renewcommand{\arraystretch}{1.3}
\begin{tabularx}{\textwidth}{
    >{\bfseries}p{3cm}
    X
    X
}    
\toprule
&
\textbf{Descripción} &
\textbf{Formato / Especificación} \\
\midrule
Entrada &
  Flujo de video desde webcam USB o cámara integrada &
  Resolución variable (640$\times$480 -- 1920$\times$1080), 30~FPS \\
\midrule
Salida &
  Frame de imagen normalizado en espacio RGB &
  NumPy array RGB, resolución fija 640$\times$480, píxeles en $[0,\,255]$ \\
\bottomrule
\end{tabularx}
\caption{Entradas y salidas del Módulo 1.}
\label{tab:m1_io}
\end{table}

\paragraph{Herramientas Tecnológicas}
\begin{itemize}[leftmargin=1.5em]
  \item \textbf{Software:} Python~3.10+ como lenguaje de implementación.
  \item \textbf{Librería principal:} OpenCV (cv2~v4.8+) \cite{Bradski2000}
        — captura de video (cv2.VideoCapture), redimensionado
        (cv2.resize) y conversión de espacio de color
        (cv2.cvtColor).
  \item \textbf{Librería de soporte:} NumPy (v1.24+) para manejo
        y operaciones sobre arrays de imagen.
\end{itemize}

\paragraph{Técnicas de Visión Computacional}
\begin{itemize}[leftmargin=1.5em]
  \item \textbf{Realzado y restauración:} aplicación opcional de filtro
        gaussiano (cv2.GaussianBlur) para atenuación de ruido
        de alta frecuencia antes de la inferencia de
        MediaPipe~\cite{GonzalezWoods2018}.
  \item Redimensionado con interpolación bicúbica
        (cv2.INTER\_CUBIC) para preservar calidad de bordes
        al escalar el frame.
  \item Conversión de espacio de color BGR$\to$RGB requerida por el
        modelo BlazePose de MediaPipe.
\end{itemize}

\subsection{Módulo 2 -- Detección de Pose con MediaPipe}
\subsubsection*{Objetivo}

Detectar los 33~landmarks del cuerpo humano en cada frame de video
utilizando el modelo MediaPipe~Pose (BlazePose), con énfasis en los
puntos clave del brazo: hombros (landmarks~11,~12), codos
(13,~14) y muñecas (15,~16)~\cite{Bazarevskyetal2020}.

\subsubsection*{Entradas y Salidas}

\begin{table}[h!]
\centering
\renewcommand{\arraystretch}{1.3}
\begin{tabularx}{\textwidth}{
    >{\bfseries}p{3cm}
    X
    X
}    
\toprule
&
\textbf{Descripción} &
\textbf{Formato / Especificación} \\
\midrule
Entrada &
  Frame RGB preprocesado (salida de M1) &
  NumPy array RGB 640$\times$480 \\
\midrule
Salida &
  33~landmarks corporales con visibilidad &
  \texttt{NormalizedLandmarkList}: coord.\
  $(x,\,y,\,z) \in [0,1]$ + visibilidad $\in [0,1]$
  por punto \\
\bottomrule
\end{tabularx}
\caption{Entradas y salidas del Módulo 2.}
\label{tab:m2_io}
\end{table}
 
\paragraph{Herramientas Tecnológicas}
\begin{itemize}[leftmargin=1.5em]
  \item \textbf{Modelo de IA preentrenado:}
        MediaPipe~Pose (BlazePose)~v0.10+~\cite{Bazarevskyetal2020,Lugaresietal2019}
        — modelo de estimación de pose de 33~landmarks entrenado sobre
        el dataset COCO~Keypoints. No requiere GPU.
  \item \textbf{Framework:} MediaPipe (Google) \cite{Lugaresietal2019}
        — framework de construcción de pipelines de percepción con
        soporte cross-platform.
  \item \textbf{Librería de soporte:} NumPy para extracción y
        transformación de las coordenadas de landmarks al espacio de
        píxeles.
\end{itemize}
 
\paragraph{Técnicas de Visión Computacional}
\begin{itemize}[leftmargin=1.5em]
  \item \textbf{Pipeline BlazePose de dos etapas:}
        \begin{enumerate}[leftmargin=2em]
          \item Detector de persona que delimita el bounding
                box de la región de interés.
          \item Estimador de pose que infiere los 33 landmarks con
                alta precisión dentro de esa región.
        \end{enumerate}
  \item \textbf{Modo de video} (static\_image\_mode=False):
        el modelo aprovecha la coherencia temporal entre frames
        consecutivos para acelerar la inferencia y estabilizar los
        landmarks.
  \item \textbf{Extracción de landmarks de interés:} se seleccionan los
        índices 11–16 (hombros, codos, muñecas) de los 33~puntos
        detectados para alimentar el módulo de cálculo angular.
\end{itemize}

\subsection{Módulo 3 -- Cálculo de Ángulos Articulares}
\subsubsection*{Objetivo}
Calcular los ángulos de flexión/extensión en las articulaciones
hombro, codo y muñeca a partir de las coordenadas 2D de los landmarks,
usando la definición vectorial del ángulo entre dos segmentos
articulares.

\subsubsection*{Entradas y Salidas}

\begin{table}[h!]
\centering
\renewcommand{\arraystretch}{1.3}
\begin{tabularx}{\textwidth}{
    >{\bfseries}p{3cm}
    X
    X
}    
\toprule
&
\textbf{Descripción} &
\textbf{Formato / Especificación} \\
\midrule
Entrada &
  Coordenadas 2D de hombro, codo y muñeca (salida de M2) &
  Pares $(x,\,y)$ en píxeles para cada articulación
  (landmarks 11–16) \\
\midrule
Salida &
  Tripleta de ángulos articulares suavizados &
  $\theta_{\text{hombro}},\,\theta_{\text{codo}},\,
   \theta_{\text{muñeca}}$ en grados $[0^\circ,\,180^\circ]$ (float64) \\
\bottomrule
\end{tabularx}
\caption{Entradas y salidas del Módulo 3.}
\label{tab:m3_io}
\end{table}

\paragraph{Herramientas Tecnológicas}\mbox{}

\begin{itemize}[leftmargin=1.5em]
  \item \textbf{Librería principal:} NumPy (v1.24+) para
        cálculo del producto escalar (np.dot), norma de
        vectores (np.linalg.norm) y función
        arccos (np.arccos) necesarios en la fórmula
        angular.
  \item \textbf{Librería de soporte:} Matplotlib (v3.7+)
        para graficar la evolución temporal de los ángulos articulares
        durante las pruebas y análisis de resultados.
\end{itemize}
 
\paragraph{Técnicas de Visión Computacional y Geometría}
\begin{itemize}[leftmargin=1.5em]
  \item \textbf{Cálculo de ángulos mediante producto escalar:}
        \cite{GonzalezWoods2018}
 
        \begin{equation}
            \theta \;=\; \arccos{ \frac{\vec{v}_1 \cdot \vec{v}_2}{|\vec{v}_1|\,|\vec{v}_2|} }
            \label{eq:angulo}
        \end{equation}

        donde $\vec{v}_1$ y $\vec{v}_2$ son los vectores que forman
        cada par de segmentos en la articulación.

  \item \textbf{Filtrado temporal con media móvil} de ventana $w = 5$
        frames para suavizar fluctuaciones por ruido de detección de
        BlazePose y evitar movimientos bruscos en la simulación del
        brazo robótico.
  \item \textbf{Cálculo de vectores entre articulaciones:} diferencia
        de coordenadas entre el landmark proximal y el distal para
        construir cada vector de segmento corporal.
\end{itemize}
 

\subsection{Módulo 4 -- Simulación del Brazo Robótico Virtual}

\subsubsection*{Objetivo}
Controlar un UR5 con gripper RG2 dentro de CoppeliaSim a partir de los
ángulos y la apertura calculados por M3, manteniendo límites operativos y
un estado de seguridad verificable~\cite{Altayeb2023}.

\subsubsection*{Entradas y Salidas}

\begin{table}[h!]
\centering
\renewcommand{\arraystretch}{1.3}
\begin{tabularx}{\textwidth}{
    >{\bfseries}p{3cm}
    X
    X
}    
\toprule
&
\textbf{Descripción} &
\textbf{Formato / Especificación} \\
\midrule
Entrada &
  Comando humano tipado (salida de M3) &
  Ángulos de hombro y codo en grados, apertura normalizada y marca temporal \\
\midrule
Salida &
  Estado y telemetría del UR5/RG2 &
  Objetivos articulares, señal del gripper, readback y snapshot de seguridad \\
\bottomrule
\end{tabularx}
\caption{Entradas y salidas del Módulo 4.}
\label{tab:m4_io}
\end{table}
 
\paragraph{Herramientas Tecnológicas}
\begin{itemize}[leftmargin=1.5em]
  \item \textbf{Simulador:} CoppeliaSim 4.10 con una escena UR5/RG2 y
        Vision Sensor.
  \item \textbf{Comunicación:} ZeroMQ Remote API para consultar la escena,
        enviar objetivos y leer el estado del robot.
  \item \textbf{Concurrencia:} worker exclusivo para las RPC y una cola lógica
        \textit{latest-only} para descartar comandos obsoletos.
\end{itemize}

\paragraph{Técnicas de Robótica y Seguridad}
\begin{itemize}[leftmargin=1.5em]
  \item Mapeo lineal de los ángulos humanos a los límites efectivos de hombro
        y codo obtenidos de la configuración y de la escena.
  \item Preflight de jerarquía, joints, RG2, intervalos, modo dinámico,
        colecciones de colisión y simulación activa.
  \item Hold, watchdog de pose, ESTOP enclavado, detección reactiva de
        colisiones y retorno verificado a \textit{home}.
  \item Orientación del modelo y encuadre del Vision Sensor durante el
        preflight para presentar el movimiento en un plano vertical estable.
\end{itemize}

\subsection{Módulo 5 -- Visualización y Métricas en Tiempo Real}

\subsubsection*{Objetivo}
Integrar en un dashboard la cámara humana anotada, la vista del Vision Sensor,
las mediciones, los diagnósticos y los controles de seguridad del sistema.

\subsubsection*{Entradas y Salidas}

\begin{table}[h!]
\centering
\renewcommand{\arraystretch}{1.3}
\begin{tabularx}{\textwidth}{
    >{\bfseries}p{3cm}
    X
    X
}    
\toprule
&
\textbf{Descripción} &
\textbf{Formato / Especificación} \\
\midrule
Entrada &
  Frames y snapshots de M1--M4 &
  Cámara humana, Vision Sensor, mediciones y estado de seguridad \\
\midrule
Salida &
  Dashboard interactivo en tiempo real &
  Dos paneles visuales, telemetría, diagnósticos y acciones de teclado \\
\bottomrule
\end{tabularx}
\caption{Entradas y salidas del Módulo 5.}
\label{tab:m5_io}
\end{table}
 
\paragraph{Herramientas Tecnológicas}
\begin{itemize}[leftmargin=1.5em]
  \item \textbf{Interfaz:} Tkinter/ttk para componer el dashboard y procesar
        las acciones del usuario sin bloquear el hilo de visión.
  \item \textbf{OpenCV:} anotación y conversión de los frames de cámara antes
        de mostrarlos en la interfaz~\cite{Bradski2000}.
  \item \textbf{MediaPipe Drawing Utils}~\cite{Lugaresietal2019}
        — dibujo automático del skeleton de 33~landmarks y
        conexiones corporales sobre el frame.
  \item \textbf{Pillow:} conversión de los frames BGR/RGB a imágenes aptas
        para los dos paneles de Tkinter.
\end{itemize}
 
\paragraph{Técnicas de Visión Computacional}
\begin{itemize}[leftmargin=1.5em]
  \item \textbf{Cálculo de FPS en tiempo real:} medición del tiempo
        de procesamiento por frame con timestamps de Python
        (time.time()) para monitoreo continuo del rendimiento
        del sistema.
  \item \textbf{Visualización compuesta:} publicación de snapshots inmutables
        con cámara humana, Vision Sensor, telemetría y diagnósticos, consumidos
        por la GUI mediante \texttt{after()}.
\end{itemize}

\section{Soluciones existentes}

Hoy en día, replicar los movimientos del cuerpo humano mediante visión por computadora es un campo clave para disciplinas como la robótica, la interacción persona-ordenador, la medicina de rehabilitación y la automatización. En la literatura actual encontramos diferentes estrategias para resolver esto, pero casi todas comparten el mismo fin: capturar lo que hace un usuario, calcular dónde están sus articulaciones en el espacio y trasladar esa información a un robot o a un entorno digital.

Una de las líneas de investigación con más fuerza es la teleoperación por imitación. Aquí el reto principal no es solo registrar el movimiento de los brazos o piernas del usuario, sino ver cómo adaptar esos datos a la estructura de un robot, cuyas articulaciones y límites físicos suelen ser muy diferentes a los de un humano.

Por ejemplo, Zhang y su equipo diseñaron un sistema para conectar los movimientos de un brazo humano con un robot industrial. Lo interesante de su propuesta es que crearon un modelo cinemático capaz de traducir las trayectorias en tiempo real; el robot lograba replicar la orientación del brazo incluso teniendo una anatomía completamente distinta \cite{wang2020fast}. Esto demostró que la compatibilidad física no es una barrera insalvable. En una línea parecida, Xuan y Tang apostaron por la simplicidad del hardware al usar una sola cámara web común (óptica monocular) para controlar un brazo robótico \cite{xuan2017vision}. Su trabajo dejó claro que no hacen falta sensores pegados al cuerpo ni trajes especiales si se cuenta con un buen procesamiento visual.

Por otro lado, cuando se trabaja con robots humanoides, la complejidad aumenta. En otro estudio, Han et al. se enfocaron en este problema desarrollando un método para rastrear específicamente el codo y la muñeca del operario. Su sistema ajusta esas coordenadas sobre la marcha para que el robot aprenda tareas de manipulación sin perder el equilibrio ni la fluidez \cite{han2024robust}. Más recientemente, un estudio de la Universidad de los Andes (2026) propuso usar cámaras RGB-D (que miden profundidad) junto con algoritmos de estimación de pose. Su investigación confirmó que la visión artificial actual es lo suficientemente madura como para interactuar con gemelos virtuales y robots físicos con una precisión bastante aceptable, dejando atrás los incómodos sensores portables \cite{alibeigi2017inverse}.

Al revisar todas estas soluciones, salta a la vista que casi todos los autores siguen una misma receta: capturan el video, detectan los puntos clave del cuerpo, analizan la geometría del movimiento y lo envían al dispositivo. Sin embargo, hay un problema: casi toda la investigación actual se empeña en usar robots físicos caros o equipos muy específicos. Es justo ahí donde este proyecto quiere marcar la diferencia. En lugar de lidiar con hardware costoso, proponemos crear una simulación con un brazo robótico virtual. De esta forma, reducimos los costos y la complejidad técnica a cero, pero manteniendo el foco en lo importante: experimentar y validar algoritmos de visión artificial y procesamiento de imágenes en tiempo real.
